\begin{proposition}[Tangent cones along the determinant rank strata]
\label{prop:determinant-tangent-strata}
Let $p\in C_m$ have complex rank $q\leq m-2$, and put $r=m-q$.
The rescalings of $C_m$ about $p$ have the unique tangent cone
\begin{equation}\label{eq:determinant-rank-tangent}
 \R^{2(m^2-r^2)}\times C_r,
 \qquad C_r=\{\operatorname{Re}\det W=0:W\in M_r(\mathbb C)\},
\end{equation}
up to an orthogonal transformation. The corresponding oriented
boundary currents and their associated varifolds also converge,
with multiplicity one. The first nonzero Taylor term of $F_m$ at
$p$ has degree $r$.
\end{proposition}

\begin{proof}
The assertion at $p=0$ is homogeneity. Otherwise, unitary left and
right changes of coordinates put $p$ in the form
$\operatorname{diag}(A_0,0)$, where $A_0$ is a positive diagonal
$q$-by-$q$ matrix. These changes are Euclidean isometries; they
replace $F_m$ by $\operatorname{Re}(\chi\det Z)$ for some
$|\chi|=1$. Write nearby matrices as
\[
 Z=\begin{pmatrix}A&B\\ C&D\end{pmatrix},\qquad
 V=D-CA^{-1}B.
\]
Block elimination, with determinant one, gives
$\det Z=(\det A)\det V$. Put
\[
 \zeta(A)=\frac{\chi\det A}{|\det A|},\qquad
 U(A)=\operatorname{diag}(\zeta(A),1,\ldots,1),\qquad
 W=U(A)V.
\]
Then
\[
 \operatorname{Re}(\chi\det Z)=|\det A|\operatorname{Re}\det W.
\]
Thus the smooth local coordinate map \(\Psi(Z)=(A-A_0,B,C,W)\)
identifies both the zero set and its negative phase with a product
of a smooth coordinate space and the corresponding determinant
zero set or phase. At $p$, the differential of the Schur complement
is $dV=dD$, and the derivative of $U(A)$ contributes nothing
because $V(p)=0$. Consequently $D\Psi_p$ is the identity on the
free blocks and the unitary row transformation $D\mapsto U(A_0)D$
on the last block. In particular, it is an isometry.

For $s\downarrow0$, the maps $s^{-1}\Psi(p+sH)$ converge in
$C^1$ on compact sets to $D\Psi_p$. Since the target zero set and
phase are products of cones, this proves local Hausdorff and
local $L^1$ convergence of their inverse images. The same local
diffeomorphisms push forward the multiplicity-one boundary
current. Their $C^1$ convergence and the area formula give local
current and mass convergence, hence convergence of the associated
varifolds. This also proves uniqueness. Finally, the determinant
identity gives the leading term
\[
 \operatorname{Re}\bigl(\chi\det(A_0)\det H_{22}\bigr),
\]
which has degree $r$ and the zero set in
\eqref{eq:determinant-rank-tangent}.
\end{proof}

\begin{example}[Cubic transverse tangents at nonvertex points]
\label{ex:cubic-nonvertex-tangents}
For $F_4=\operatorname{Re}\det Z$ on $M_4(\mathbb C)$, take
$p=\operatorname{diag}(1,0,0,0)$. Then
\[
 \operatorname{Tan}_p C_4=\R^{14}\times
 \{\operatorname{Re}\det W=0:W\in M_3(\mathbb C)\}.
\]
The transverse equation is cubic. In contrast, the rank-two stratum
has quadratic transverse tangents. Both the original cone and the
cubic factor are area minimizing by \cref{main:complex-determinants}.
Their regular minimality is the determinant construction of
\cite[Corollary~5.2 and Example~5.4]{HoppeTkachev19}.

The normed-product family gives an equally direct example. Let
$P_d(x_1,\ldots,x_d)$ denote the left-associated octonionic product,
and let $\mathcal O_d=\{\operatorname{Re}P_d=0\}$. At
$p=(1,0,0,0)\in\mathcal O_4$, the coordinate change
\[
 (a,b,c,d)\longmapsto(a,(1+a)b,c,d)
\]
is a local diffeomorphism whose derivative at zero is the identity.
It transforms $\operatorname{Re}P_4(1+a,b,c,d)$ into
$\operatorname{Re}((b'c)d)$. Hence the unique set and boundary-current
tangent is $\R^8\times\mathcal O_3$. The same argument gives
$\R^8\times\mathcal O_r$ at $(1,0,\ldots,0)\in\mathcal O_{r+1}$
for every $r\geq2$. These cones are irreducible and area minimizing
by \cref{prop:norm-product-closure,thm:octonionic-product-minimizing}.

These are not cylindrical extensions of their defining equations.
For determinants, evaluation of directional derivatives on
rank-$(m-1)$ matrices detects every matrix entry of the direction.
For normed products, the directional derivative terms have distinct
block multidegrees, and $\operatorname{Re}(uv)$ is a nondegenerate
bilinear form. In both cases $D_vF\equiv0$ forces $v=0$.
The Euclidean factor appears in the tangent at the singular point;
the radial line splitting also follows from
\cref{sec:tangent-instability-transfer}.
\end{example}


For comparison, if a homogeneous cubic $F$ is defined on its
essential space, every nonzero critical point $p$ has
$\Hess F(p)\ne0$, since $D_pF(x)=\tfrac12\Hess F(p)[x,x]$
and essentiality excludes a nonzero direction with $D_pF\equiv0$.
Thus its first nonzero Taylor term is quadratic; for the harmonic
minimal case, see \cref{prop:cubic-general-singular-spectrum,sec:tangent-instability-transfer}.
The higher-order transverse cones above occur in defining degree four and higher.
